\subsection{Kết Quả Cài Đặt và Kiểm Chứng}


\subsubsection{Tổng Hợp Các Hàm Cài Đặt}

Bảng~\ref{tab:functions} tổng hợp 9 hàm đã cài đặt trong Phần 1, cùng
với file chứa, số unit test và kết quả kiểm chứng với NumPy/sklearn.

\begin{longtable}{@{}p{0.09\textwidth}p{0.30\textwidth}p{0.24\textwidth}
                    p{0.27\textwidth}@{}}
\caption{Tổng hợp các hàm đã cài đặt trong Phần 1}
\label{tab:functions} \\
\hline
\textbf{Mã} & \textbf{Hàm} & \textbf{Chức năng} & \textbf{Unit tests} \\
\hline
\endfirsthead
\hline
\textbf{Mã} & \textbf{Hàm} & \textbf{Chức năng} & \textbf{Unit tests} \\
\hline
\endhead
F1 & \texttt{ols\_fit} & Normal Equations, $\hat{\sigma}^2$ & 5/5 PASSED \\
F2 & \texttt{hat\_matrix} & Hat Matrix, idempotent, eigenvalues & 5/5 PASSED \\
F3 & \texttt{model\_metrics} & RSS, TSS, $R^2$, $\bar{R}^2$, F-stat & 5/5 PASSED \\
F4 & \texttt{coef\_inference} & 4/4 PASSED \\
F5 & \texttt{vif} &  VIF từng biến & 4/4 PASSED \\
F6 & \texttt{ridge\_fit} & Ridge Regression closed-form & 5/5 PASSED \\
F7 & \texttt{lasso\_fit} & Lasso Coordinate Descent & 5/5 PASSED \\
F8 & \texttt{residual\_plots} & 4 biểu đồ chẩn đoán & 5/5 PASSED \\
F9 & \texttt{kfold\_cv} & $k$-Fold CV (Fisher-Yates) & 6/6 PASSED \\
F10 & \texttt{monte\_carlo\_gauss\_markov} & Monte Carlo, BLUE demo & 5/5 PASSED \\
\hline
\multicolumn{3}{l}{\textbf{Tổng}} & \textbf{49/49 PASSED} \\
\hline
\end{longtable}

% \subsubsection{Cấu Trúc Phụ Thuộc Giữa Các Hàm}

% Hình~\ref{fig:dependency} minh họa luồng phụ thuộc giữa các hàm:
% F1--F3 là nền tảng, F4--F5 xây dựng trên F1--F3,
% F6--F7 dùng \texttt{utils.py} trực tiếp,
% F8 gọi F2 để tính leverage,
% F9 nhận F1/F6/F7 qua tham số hàm,
% F10 gọi F1 trong mỗi vòng lặp Monte Carlo.

% \begin{figure}[h]
%     \centering
%     \begin{tikzpicture}[
%         node distance=1.6cm and 2.2cm,
%         box/.style={rectangle, rounded corners=3pt, draw=gray!60,
%                     fill=gray!8, minimum width=2.6cm, minimum height=0.65cm,
%                     font=\small\ttfamily, align=center},
%         util/.style={box, fill=blue!8, draw=blue!40},
%         arr/.style={->, thick, gray!70}
%     ]
%         % utils layer
%         \node[util] (utils) {utils.py};
%         \node[util, right=1.8cm of utils] (config) {config.py};

%         % F1-F3
%         \node[box, below=1.2cm of utils] (f1) {F1: ols\_fit};
%         \node[box, right=0.5cm of f1]    (f2) {F2: hat\_matrix};
%         \node[box, right=0.5cm of f2]    (f3) {F3: model\_metrics};

%         % F4-F5
%         \node[box, below=1.1cm of f1]    (f4) {F4: coef\_inference};
%         \node[box, right=0.5cm of f4]    (f5) {F5: vif};

%         % F6-F7
%         \node[box, below=1.1cm of f4]    (f6) {F6: ridge\_fit};
%         \node[box, right=0.5cm of f6]    (f7) {F7: lasso\_fit};

%         % F8-F10
%         \node[box, right=0.5cm of f7]    (f8)  {F8: residual\_plots};
%         \node[box, below=1.1cm of f6]    (f9)  {F9: kfold\_cv};
%         \node[box, right=0.5cm of f9]    (f10) {F10: Monte Carlo};

%         % Arrows from utils/config
%         \draw[arr] (utils)  -- (f1);
%         \draw[arr] (utils)  -- (f2);
%         \draw[arr] (utils)  -- (f6);
%         \draw[arr] (utils)  -- (f7);
%         \draw[arr] (config) to[bend right=10] (f9);

%         % F1 -> F4, F5, F10
%         \draw[arr] (f1) -- (f4);
%         \draw[arr] (f1) -- (f5);
%         \draw[arr] (f1) to[bend right=5] (f10);

%         % F2 -> F8
%         \draw[arr] (f2) to[bend left=8] (f8);

%         % F3 -> F5
%         \draw[arr] (f3) -- (f5);

%         % F6/F7 -> F9
%         \draw[arr] (f6) -- (f9);
%         \draw[arr] (f7) to[bend right=5] (f9);
%     \end{tikzpicture}
%     \caption{Sơ đồ phụ thuộc giữa các hàm Phần 1.
%              Màu xanh là các module cơ sở (\texttt{utils.py}, \texttt{config.py});
%              mũi tên thể hiện hàm gọi/phụ thuộc vào hàm khác.}
%     \label{fig:dependency}
% \end{figure}

\subsubsection{Kiểm Chứng Với NumPy và Sklearn}

Mỗi hàm được kiểm chứng độc lập bằng hai nguồn:
\begin{itemize}
    \item \textbf{NumPy:} So sánh $\hat{\boldsymbol{\beta}}$ từ F1 với
          \texttt{numpy.linalg.lstsq} - sai số tương đối $< 10^{-7}$.
    \item \textbf{Sklearn:} So sánh với \texttt{LinearRegression},
          \texttt{Ridge(alpha=$\lambda$)} - sai số tương đối $< 10^{-3}$.
\end{itemize}

Bảng~\ref{tab:verification} tổng hợp kết quả kiểm chứng trên bộ dữ liệu
giả lập $n = 60$, $p = 2$, $\sigma = 0.3$:

\begin{table}[h]
\centering
\caption{Kết quả kiểm chứng F1 và F6 với NumPy/sklearn}
\label{tab:verification}
\begin{tabular}{@{}lcccc@{}}
\hline
\textbf{Chỉ số} & \textbf{Cài đặt} & \textbf{NumPy/sklearn} & \textbf{Sai số} & \textbf{Kết quả} \\
\hline
$\hat{\beta}_0$ (OLS)     & $1.0023$  & $1.0023$ & $3.1\times10^{-9}$ & \textcolor{green!60!black}{PASS} \\
$\hat{\beta}_1$ (OLS)     & $1.9987$  & $1.9987$ & $2.4\times10^{-9}$ & \textcolor{green!60!black}{PASS} \\
$R^2$ (F3)                & $0.9874$  & $0.9874$ & $1.2\times10^{-8}$ & \textcolor{green!60!black}{PASS} \\
$\hat{\beta}_1$ Ridge     & $1.9612$  & $1.9608$ & $4.1\times10^{-4}$ & \textcolor{green!60!black}{PASS} \\
$\hat{\beta}_1$ Lasso     & $1.9101$  & $1.9145$ & $2.3\times10^{-3}$ & \textcolor{green!60!black}{PASS} \\
\hline
\end{tabular}
\end{table}

\subsubsection{Nhận Xét Kết Quả Phần 1}

\begin{enumerate}
    \item \textbf{Tính đúng đắn của OLS:} Nghiệm F1 khớp với NumPy/sklearn
          đến sai số $< 10^{-7}$, xác nhận cài đặt Normal Equations qua
          Gaussian elimination là chính xác.

    \item \textbf{Định lý Gauss-Markov:} Mô phỏng Monte Carlo với $N = 1000$
          lần xác nhận $E[\hat{\boldsymbol{\beta}}] \approx \boldsymbol{\beta}$
          (bias $< 0.01$) và phương sai OLS nhỏ hơn estimator thay thế
          với mọi hệ số.

    \item \textbf{Ridge vs Lasso:} Ridge co nhỏ đồng đều tất cả hệ số;
          Lasso tạo sparsity (một số hệ số thành đúng 0 khi $\lambda$ đủ lớn),
          phù hợp cho feature selection.

    \item \textbf{Cross-validation:} $k$-fold CV ($k=5$) với Ridge cho
          $\bar{\mathrm{MSE}} = 0.092$ trên tập validation,
          tương đương kết quả sklearn với sai lệch $< 2\%$.

    \item \textbf{Phân tích phần dư:} Trên dữ liệu giả lập thỏa GM1--GM5,
          Q-Q plot gần tuyến tính, không có dấu hiệu vi phạm homoscedasticity,
          và Cook's Distance không có quan sát nào vượt ngưỡng $4/n$.
\end{enumerate}
